\section{Introduction}
\label{sec:intro}

In the quest to create versatile, multi-task artificial intelligence, model merging has emerged as a compelling and computationally efficient alternative to joint multi-task training \cite{wortsman2022model, chen2024survey}. Rather than retraining billions of parameters from scratch—which demands astronomical compute budgets and private, centralized datasets—model merging directly fuses task-specific expert checkpoints that have been fine-tuned from a shared pretrained base model \cite{ilharco2023editing, yadav2024ties}. Yet, despite its practical appeal, the foundational mathematics governing modern model merging remain surprisingly primitive. Virtually all existing methods are constrained by Euclidean linearity: they rely on static, coordinate-wise operations such as weighted averages, linear interpolations, or heuristic sign-agreement masks in flat parameter space \cite{yadav2024ties, yu2023language}. 

This flat, Euclidean perspective ignores a fundamental truth of deep neural networks: the loss landscape of high-dimensional parameter spaces is highly non-convex, warped, and characterized by intricate topological structures \cite{li2018visualizing, fort2019deep}. Performing simple algebraic average updates between task experts is akin to drawing a straight line through a mountainous terrain; the resulting path frequently traverses high-loss barriers and saddle points, destroying task-specific representations and inducing severe parameter interference \cite{yang2024adamerging, zhang2024task}. 

In this work, we propose a fresh, intuitive perspective on parameter merging and test-time adaptation. To find deep connections between physical phenomena and deep learning, we frame parameter trajectories as an intuitive physical analog: \emph{What if neural network parameters were treated as continuous physical fluid phases flowing and coalescing through a viscous vector field?}

To realize this perspective, we introduce \textbf{FluidMerge} (Fluid-Dynamic Parameter Coalescence), a framework that interprets the test-time adaptation of merged models as a continuous-time advection-diffusion fluid flow. Drawing inspiration from classical fluid-dynamic principles, we treat the trajectory of parameters $\theta(t)$ over a virtual time horizon $t \in [0, T]$ (where $T > 0$) as a mobile fluid mixture. Under this analog, the parameter trajectory is governed by two cooperating forces:
\begin{enumerate}
    \item \textbf{Advection (Task-Gradient Pull):} Self-supervised task-specific losses generate a velocity field that pulls the weight particles along low-loss streamlines towards the task-specific basins.
    \item \textbf{Diffusion (Viscous Parameter Regularization):} Parameter diffusion acts as physical fluid viscosity to smooth out updating velocities and prevent representation tearing. While we initially introduce a grid-based spatial Laplacian operator over weight coordinates as an intuitive baseline metaphor, we show that this grid-based approach commits a fundamental representation category error due to neural permutation invariance. To resolve this, we place the primary emphasis of the parameter-space "fluid flow" on the Riemannian manifold of the Fisher metric from the outset, formulating a mathematically sound, coordinate-free, and permutation-invariant Fisher viscosity.
\end{enumerate}

Rather than framing this fluid analogy as an entirely new set of mathematical primitives, we ground it rigorously in established machine learning foundations. We solve the continuous trajectory $\theta(t)$ numerically using a discretized Euler integration scheme, guided by teacher-provided soft labels on unlabeled test-time data.

Our rigorous empirical evaluation of FluidMerge alongside state-of-the-art baselines—such as SyMerge \cite{jung2025symerge}, AdaMerging \cite{yang2024adamerging}, and Task Surgery \cite{zhang2024task}—across eight diverse classification tasks using a Vision Transformer (ViT-B-32) \cite{radford2021clip, dosovitskiy2020image} yields profound and unexpected insights into the nature of deep representations. Specifically, our experiments expose a deep-seated \textbf{domain shift barrier} between pretrained base models and task-specific classifiers. Because the classification heads are trained on top of fully fine-tuned expert checkpoints, they are highly sensitive to the exact coordinate shifted representations of their respective experts. When paired with an unaligned base image encoder, they output pure noise, yielding random-guess performance (averaging 6.23\% top-1 accuracy). 

Furthermore, our analysis of the test-time adaptation dynamics reveals a critical failure mode: under unaligned encoder features, self-supervised test-time optimization (both in FluidMerge and the baselines) causes the classifier heads to rapidly overfit to pseudo-labels. Consequently, the Expected Calibration Error (ECE) explodes to over 90%—the model becomes highly confident in its incorrect predictions on unseen validation data. This analysis exposes a critical representational domain shift bottleneck, highlighting that post-hoc adaptation on unlabeled test data struggles to bridge the deep representation gap from scratch without prior structural weight initializations.

To resolve these core challenges, we design, implement, and validate two key advancements within our framework. Importantly, we de-escalate metaphorical overselling by explicitly identifying their exact mathematical equivalences to standard deep learning techniques:
\begin{enumerate}
    \item \textbf{Expert-Weighted Initial Boundary Conditions (Task Arithmetic Initialization):} Instead of initializing our fluid simulation at the pure pretrained base weights $\theta(0) = \theta_0$, we initialize it at a merged weight state incorporating task-expert vectors. Mathematically, this is identical to initializing the optimization starting from the standard \textbf{Task Arithmetic} average weight-vector ($\theta_{\text{TA}}$) of \citet{ilharco2023editing}. This places the parameters within high-performing multi-task basins from the outset, using the continuous-time flow to smoothly refine and align representation boundaries rather than attempting to reconstruct them from scratch.
    \item \textbf{Fisher-Information-based Viscosity (Continuous Elastic Weight Consolidation):} We replace the coordinate-dependent spatial Laplacian with a mathematically sound, coordinate-free, and permutation-invariant viscosity operator based on the empirical diagonal Fisher Information Matrix. Rigorous mathematical deconstruction reveals that under discretized Euler integration, this viscosity operator is \textbf{mathematically isomorphic to running gradient descent under an Elastic Weight Consolidation (EWC) penalty} \cite{kirkpatrick2017overcoming} that anchors parameters to their initial Task Arithmetic state. In this formulation, "viscosity" acts as a coordinate-wise restorative spring force that stabilizes high-information parameters, preventing functional tearing and representation collapse.
\end{enumerate}

Our rigorous empirical evaluation on the Vision Transformer (ViT-B-32) across eight classification datasets demonstrates that these advancements successfully unlock the power of weight-fluid simulations. Under the Expert-Weighted Initial Boundary Conditions, FluidMerge achieves outstanding performance, successfully adapting representation boundaries and resolving the domain shift barrier, while Fisher-Information viscosity successfully stabilizes the continuous-time trajectory, preventing calibration collapse and representation tearing, outperforming static Task Arithmetic, head-tuning controls, and standard $L_2$ weight anchoring baselines.

To summarize, our primary contributions are:
\begin{itemize}
    \item \textbf{The FluidMerge Perspective:} We propose an intuitive and elegant perspective framing model merging as a continuous-time physical fluid flow, bridging deep learning parameter trajectories with fluid mechanics.
    \item \textbf{Fisher-Information Viscosity \& EWC Isomorphism:} We introduce a mathematically sound, coordinate-free, and permutation-invariant parameter-space viscosity regularizer based on the empirical Fisher metric, proving its exact mathematical isomorphism to a continuous-time formulation of Elastic Weight Consolidation.
    \item \textbf{Expert-Weighted Initial Boundary Conditions:} We formulate and implement boundary conditions that initialize the parameter fluid at a merged task-expert state, mathematically equivalent to Task Arithmetic initialization, completely bypassing the representation reconstruction bottleneck.
    \item \textbf{Empirical Resolution of the Domain Shift Barrier:} We systematically demonstrate that combining Task Arithmetic initialization with Fisher-Information-based continuous-time advection-diffusion flow successfully adapts high-performing representations, outperforming standard static weight merging and standard $L_2$ anchoring.
\end{itemize}
